\chapter{可测函数}
\label{chap:measurable-functions}

连续函数用开集原像保持开集来表达。测度空间没有距离和邻域，能够保留的
条件是：目标空间中可测集合的原像仍可测。这个定义足以控制函数的水平集、
逐点极限和简单函数逼近，也是 Lebesgue 积分的入口。

\section{可测映射与半直线判据}

\begin{definition}
设 $(X,\mathcal A)$、$(Y,\mathcal B)$ 为可测空间。映射 $f:X\to Y$ 称为
\term{可测}，如果对每个 $B\in\mathcal B$，都有
$f^{-1}(B)\in\mathcal A$。
\end{definition}

连续映射在拓扑空间的 Borel $\sigma$-代数之间可测，因为开集原像为开集，
而“原像可测”的集合族本身是 $\sigma$-代数。复合可测映射仍可测。

实值函数常允许取 $\pm\infty$，目标空间写作
$\overline\R=[-\infty,+\infty]$，采用由区间生成的 Borel
$\sigma$-代数。

\begin{theorem}[半直线判据]
\label{thm:measurable-halfline-criteria}
对 $f:X\to\overline\R$，以下条件等价：
\begin{enumerate}[label=\textup{(\roman*)}]
 \item $f$ 可测；
 \item 对每个 $a\in\R$，$\{f>a\}\in\mathcal A$；
 \item 对每个 $a\in\R$，$\{f\ge a\}\in\mathcal A$；
 \item 对每个 $a\in\R$，$\{f<a\}\in\mathcal A$；
 \item 对每个 $a\in\R$，$\{f\le a\}\in\mathcal A$。
\end{enumerate}
只检验有理数 $a$ 也足够。
\end{theorem}

\begin{proof}
可测性立即给出各水平集可测。反过来，半直线生成
$\mathcal B(\overline\R)$，故相应原像条件推出可测性。各严格与非严格
版本由
\[
 \{f>a\}=\bigcup_{n\ge1}\{f\ge a+1/n\},
 \qquad
 \{f\ge a\}=\bigcap_{n\ge1}\{f>a-1/n\}
\]
及取补集互相推出。任意实端点可用有理数单调逼近。扩充实线的端点也被
这些集合运算包括：$\{+\infty\}=\bigcap_{n\ge1}(n,+\infty]$，
$\{-\infty\}=\bigcap_{n\ge1}[-\infty,-n)$。有限端点的开区间可由半直线
取差得到，所以检验半直线没有遗漏端点或其他Borel集合。
\end{proof}

\section{代数运算与序运算}

若 $f,g$ 是有限实值可测函数，则 $f+g$、$fg$、$|f|$、
$f\vee g$、$f\wedge g$ 均可测。一个简洁证明是映射
$x\mapsto(f(x),g(x))$ 从 $X$ 到 $\R^2$ 可测：开矩形原像是两个可测
原像的交，而开集可由至多可数个有理矩形组成。再与 $\R^2$ 上连续的加法、
乘法、绝对值等复合。

对扩充实值函数，$+\infty+(-\infty)$ 没有定义，必须先排除这种点；不能
以“可测函数可相加”掩盖运算本身未定义。正部和负部
\[
 f^+=f\vee0,\qquad f^-=(-f)\vee0
\]
总是非负可测，并满足 $f=f^+-f^-$、$|f|=f^++f^-$。对每个点，
$f^+$ 与 $f^-$ 至少一个为零，所以差式不会出现 $\infty-\infty$。
$f$ 可以在不同点分别取正、负无穷；这与在同一点进行未定义的运算不同。

\begin{theorem}[可数序运算]
\label{thm:measurable-countable-order-operations}
若 $f_n:X\to\overline\R$ 可测，则
\[
 \sup_nf_n,\quad \inf_nf_n,\quad
 \limsup_nf_n,\quad\liminf_nf_n
\]
均可测。若 $f_n(x)$ 逐点收敛，则极限函数可测。
\end{theorem}

\begin{proof}
例如
\[
 \left\{\sup_nf_n>a\right\}=\bigcup_n\{f_n>a\},
 \qquad
 \inf_nf_n=-\sup_n(-f_n).
\]
再用
\[
 \limsup_nf_n=\inf_N\sup_{n\ge N}f_n,
 \qquad
 \liminf_nf_n=\sup_N\inf_{n\ge N}f_n.
\]
逐点极限存在时两者相等。
\end{proof}

\section{简单函数逼近}

\begin{definition}
只取有限多个实值且可测的函数称为\term{简单函数}。它可写成
\[
 \varphi=\sum_{k=1}^ma_k\mathbf1_{E_k},
\]
其中 $E_k$ 可测。一般表示中的 $E_k$ 可以重叠，此时一个点的函数值是
所有相关系数之和，不能把每个 $E_k$ 直接当作水平集。
若把 $a_1,\ldots,a_m$ 取为全部不同的实际函数值，再令
$E_k=\{\varphi=a_k\}$，便得到两两不交且覆盖 $X$ 的规范表示。
只列非空水平集，实际取到零时也列入零水平集；在此约定下，规范表示除
排列外唯一。若 $X=\varnothing$，规范表示为空和。
\end{definition}

\begin{theorem}[非负简单函数逼近]
\label{thm:measurable-simple-approximation}
对每个非负可测函数 $f:X\to[0,+\infty]$，存在非负简单函数列
$\varphi_n$，使
\[
 0\le\varphi_n\le\varphi_{n+1}\le f,
 \qquad \varphi_n(x)\uparrow f(x).
\]
若 $f$ 有界，可使收敛一致。
\end{theorem}

\begin{proof}
把 $[0,n)$ 按长度 $2^{-n}$ 分成小区间，定义
\[
 \varphi_n(x)=
 \begin{cases}
  k2^{-n},&k2^{-n}\le f(x)<(k+1)2^{-n},\quad
             0\le k<n2^n,\\
  n,&f(x)\ge n.
 \end{cases}
\]
这里 $n\ge1$，$k$ 为整数。水平集由半直线的差得到，故 $\varphi_n$ 简单可测。
若 $f(x)<n$，下一步网格把每个旧区间二等分，新的左端点不小于旧的左端点；
若 $f(x)\ge n$，则 $\varphi_n(x)=n$，而下一步截断高度为 $n+1$，
网格又包含 $n$，故 $\varphi_{n+1}(x)\ge n$。因此这一个序列本身已递增。
当 $f(x)<\infty$，取 $n>f(x)$ 后有
$0\le f(x)-\varphi_n(x)<2^{-n}$；当 $f(x)=+\infty$，有
$\varphi_n(x)=n\to+\infty$。两种情形都得到 $\varphi_n(x)\uparrow f(x)$。
若 $f\le M$，当 $n>M$ 时截断不再出现，误差至多 $2^{-n}$，故一致收敛。
\end{proof}

简单函数保留了有限求和的可控性，又能从下逼近任意非负可测函数。下一章
将先定义简单函数积分，再对所有下方简单函数取上确界。

\section{零测修改与可测代表}

设 $f,g$ 是同一测度空间上的扩充实值函数。若空间完备、$f$ 可测且
$g=f$ a.e.，则 $g$ 可测。事实上，对任意
$a$，两个水平集 $\{g>a\}$ 与 $\{f>a\}$ 只相差一个零测集。若空间不
完备，这个结论可能失败：取不可测的零测子集 $Z$，把零函数在 $Z$ 上改为
$1$，所得函数的水平集恰为 $Z$。

因此，后文把几乎处处相等函数视为同一元素时，将默认在完备化中选取可测
代表，或明确保留原测度空间的可测代表。

\section{Lusin 定理}

可测函数不必连续，但在有限测度的 Euclidean 集合上，可以删去任意小的
测度后得到连续限制。

\begin{theorem}[Lusin]
\label{thm:lusin}
设 $E\subset\R^d$ Lebesgue 可测且 $m(E)<\infty$，
$f:E\to\R$ 可测。对每个 $\varepsilon>0$，存在紧集 $K\subseteq E$，
使 $m(E\setminus K)<\varepsilon$，且 $f|_K$ 连续。
\end{theorem}

\begin{proof}
先删去函数值过大的部分。由于 $f$ 处处有限，集合
$E_M=E\cap\{|f|\le M\}$ 递增到 $E$；上升连续性给出某个 $M$ 使
$m(E\setminus E_M)<\varepsilon/4$。

在 $E_M$ 上先平移为非负有界函数 $f+M$，用二进网格向下取整，再减去 $M$，
得到简单函数 $\varphi_n$，使
$\norm{f-\varphi_n}_{\infty,E_M}\le2^{-n}$。把 $\varphi_n$ 的非空水平集
记为 $A_{n,1},\ldots,A_{n,r_n}$；它们构成 $E_M$ 的有限可测分割。由内
正则性，可在每个 $A_{n,j}$ 中取紧集 $K_{n,j}$，并把误差分配成
\[
 \sum_{j=1}^{r_n}m(A_{n,j}\setminus K_{n,j})
 <\varepsilon2^{-n-2}.
\]
令 $K_n=\bigcup_jK_{n,j}$。忽略其中的空集后，有限个互不相交的紧集之间
有正距离（至多一个非空项时无需距离估计），因此
$\varphi_n|_{K_n}$ 连续；同时
$m(E_M\setminus K_n)<\varepsilon2^{-n-2}$。

令
\[
 K=K_0\cap\bigcap_{n=1}^{\infty}K_n,
\]
其中先由内正则性取紧 $K_0\subseteq E_M$，使
$m(E_M\setminus K_0)<\varepsilon/4$。集合 $K$ 是紧集 $K_0$ 的闭子集，
故仍紧。并且
\[
 m(E\setminus K)
 \le m(E\setminus E_M)+m(E_M\setminus K_0)
    +\sum_{n=1}^{\infty}m(E_M\setminus K_n)
 <\varepsilon.
\]
每个 $\varphi_n|_K$ 连续，并且
$\varphi_n\to f$ 在 $K$ 上一致。因此 $f|_K$ 是连续函数的一致极限，故
连续。
\end{proof}

有限测度使“大函数值”部分可以被删去；正则性把可测水平集换成紧集；一致
逼近把简单函数的连续性传给 $f$。定理没有声称 $f$ 在某个开集上连续，也
没有把不连续点集本身变成零测集。

\section*{习题}
\addcontentsline{toc}{section}{习题}

\subsection*{A组　判据与运算}
\begin{enumerate}[label=\textbf{\thechapter.\arabic*.}]
 \exerciseitem{measurable-rational-thresholds}
 证明半直线判据只需检验有理端点。
 \exerciseitem{measurable-vector-map}
 若 $f_1,\ldots,f_n$ 为实值可测函数，证明
 $x\mapsto(f_1(x),\ldots,f_n(x))$ 是到 $\R^n$ 的 Borel 可测映射。
 \exerciseitem{measurable-algebra-operations}
 由上一题证明有限实值可测函数对加、乘、绝对值和有限最大值封闭。
 \exerciseitem{measurable-limsup-formula}
 设 $f_n:X\to\overline\R$ 为同一可测空间上的可测函数，$a\in\R$。
 写出 $\{\limsup f_n>a\}$ 的可数集合运算表达式，并据此直接证明可测性。
 \exerciseitem{measurable-convergence-set}
 设 $f_n:X\to\R$ 为同一可测空间上的可测函数。分别证明：有限实极限的
 收敛点集可测；允许极限为 $\pm\infty$ 时的收敛点集也可测。
\end{enumerate}

\subsection*{B组　简单函数与零测修改}
\begin{enumerate}[label=\textbf{\thechapter.\arabic*.},resume]
 \exerciseitem{simple-function-canonical-form}
 证明简单函数可以写成按不同取值水平集分解的规范形式，并说明该形式唯一。
 \exerciseitem{nested-dyadic-simple-approximation}
 对非负可测函数，改用截断高度 $2^n$ 构造逐项递增的二进简单函数逼近，
 并直接验证单调性及逐点收敛。
 \exerciseitem{bounded-simple-uniform-approximation}
 证明有界可测函数可由简单函数一致逼近。
 \exerciseitem{ae-modification-complete-space}
 证明完备测度空间中，可测函数的几乎处处修改仍可测；并给出非完备空间的
 反例。
 \exerciseitem{indicator-measurability}
 证明 $\mathbf1_E$ 可测当且仅当 $E$ 可测。
\end{enumerate}

\subsection*{C组　Lusin 定理及应用}
\begin{enumerate}[label=\textbf{\thechapter.\arabic*.},resume]
 \exerciseitem{lusin-simple-functions}
 先对有限值简单函数证明 Lusin 定理。
 \exerciseitem{lusin-proof-measure-budget}
 核对定理 \ref{thm:lusin} 证明中全部删去集合的测度预算，并说明所得 $K$
 为什么紧。
 \exerciseitem{lusin-continuous-extension}
 由 Tietze 延拓定理证明：Lusin 定理中的 $f|_K$ 可延拓为 $\R^d$ 上的连续
 函数。
 \exerciseitem{measurable-almost-continuous-not-continuous-ae}
 在 $[a,b]$（$a<b$）上用 Dirichlet 函数说明“对每个 $\varepsilon$ 可删去小测度后连续”不等价于
 “函数的不连续点集零测”。
 \optionalexerciseitem{measurable-functions-generated-sigma}
 对任意实值函数族 $f_i:X\to\R$（$i\in I$），证明 $X$ 上使全部
 $f_i$ 可测的最小 $\sigma$-代数是
 $\sigma\{f_i^{-1}(B):i\in I,\ B\in\mathcal B(\R)\}$。
\end{enumerate}
